\section{Elementos visuales y gestión de recursos}
\subsection{Fotogramas clave y marcas de tiempo}
Esta clase se apoya en la imagen de portada y en 4 fotogramas clave para reforzar la evidencia visual. Al no haber slides, estos fotogramas muestran conceptos centrales como self-rewarding, self-instruct y meta-judge. La tabla adjunta enumera el nombre del fotograma, el intervalo de tiempo y la descripción textual, para facilitar su consulta posterior.

\begin{table}[H]
\centering
\begin{tabular}{p{3cm}p{3cm}p{7cm}}
\toprule
Archivo del fotograma & Tiempo & Descripción \\
\midrule
frame-01-selfreward.jpg & 00:02:35--00:02:37 & Diagrama conceptual del self-improvement pipeline \\
frame-02-iterative.jpg & 00:36:00--00:36:02 & Bucle iterativo de Self-Rewarding \\
frame-03-selfinstr.jpg & 00:44:40--00:44:42 & Tareas de Self-Instruct y verified response \\
frame-04-meta.jpg & 01:01:30--01:01:32 & Meta Judge compara draft y verified response \\
\bottomrule
\end{tabular}
\caption{Fotogramas clave y marcas de tiempo usados en esta clase.}
\end{table}

\subsection{Subtítulos y apoyo para la redacción}
A partir de \texttt{lecture11.en.srt} depuramos \texttt{subs\_clean.txt}, extrayendo oraciones temáticas en bloques de 2500 líneas y dividiendo después las secciones según la lógica didáctica. La versión clean del texto reduce las repeticiones (como «days I'm using AI...»), lo que aporta material estable para ampliar el contenido de cada sección.

\subsection{Resumen de la sección}
Los recursos visuales y la depuración de subtítulos constituyen el doble sustento del material complementario: los primeros satisfacen el requisito del usuario de frame + tiempo; los segundos dan un fundamento a la redacción de las secciones y mantienen la consistencia de los detalles técnicos y del orden.

